\begin{lemma}[Two exit legs from closed noninterchangeability]
Let $G=(V,E)$ be a finite loopless multigraph whose elements of $E$ are
the exact edge identities, let $T$ be a tree-sequence, and let $\varphi$
be a proper partial $k$-edge-coloring.  Suppose that $V(T)$ is elementary
under $\varphi$.  Let $a,b\in[k]$ be distinct colors such that
$a\in\overline\varphi(T)$, the singleton $\{a\}$ is closed on $T$, and
$a,b$ are not $T$-interchangeable.

Then there are vertices $a_1,b_1,a_2,b_2$, two walks $P_1,P_2$ in the
two-color incidence graph from $a_i$ to $b_i$, and exact edges $f_1,f_2$
such that, for $i=1,2$, $P_i$ is a $T$-exit path with outside end $a_i$
and tree end $b_i$, while
\[
 a_1\ne a_2,\qquad b_1\ne b_2,\qquad
 \operatorname{supp}(P_1)\cap\operatorname{supp}(P_2)=\varnothing.
\]
Moreover, $f_i$ lies on $P_i$, is incident with $b_i$, and is an exact
$b$-colored boundary-edge identity of $V(T)$.
\end{lemma}
\begin{proof}
Put $X=V(T)$, with $X$ understood as the exact finite set defined by
\noderef{TreeSequenceVertexSet}.  Negating the universal assertion in
\noderef{TreeInterchangeable} gives nodes $z_1,z_2$ such that the two sets
\[
 C_i=\operatorname{Comp}_{a,b}(z_i)\quad(i=1,2)
\]
from \noderef{TwoColorComponent} are unequal
\noderef{TwoColorPathComponent}s, and each $C_i$ contains a vertex-node
belonging to $X$.  The equality-or-disjointness assertion of
\noderef{TwoColorComponentsDisjoint} therefore makes $C_1$ and $C_2$
disjoint as sets of incidence-graph nodes.  In particular, they cannot
meet at either a vertex-node or an exact edge-node.

We first record the precise shape of every $C_i$ that contains an edge-node.
Because it also contains a vertex-node of $X$, connectedness shows that its
edge-node is not an isolated uncolored edge-node; it is joined along the
component to a vertex-node.  Thus the path-component certificate together
with this edge-node is exactly a
\noderef{TwoColorNontrivialPathComponent}.  Apply
\noderef{TwoColorComponentSpanningPath}.  It gives a simple path $R_i$
whose support is all of $C_i$.  Properness, in the sense of
\noderef{IsProperEdgeColoring}, and
\noderef{TwoColorIncidenceDegreeAtMostTwo} bound the degree of every node
of $C_i$ by two.  Hence an internal node of $R_i$ has no neighbor outside
its two path neighbors.  The two ends cannot have an extra adjacency to
one another: that would make every node of $C_i$ have degree two, contrary
to the node with at most one neighbor supplied by
\noderef{TwoColorPathComponent}.  Nor can an end have any other extra
neighbor, since that neighbor already lies on the spanning support and
would give an internal node or that end a third neighbor.  Consequently
$C_i$ itself is the simple nontrivial path $R_i$, and the two ends of
$R_i$ have degree one in the full two-color incidence graph.

Both ends of $R_i$ are vertex-nodes.  Indeed, by
\noderef{TwoColorIncidenceGraph}, the neighbors of a colored edge-node
$f$ are precisely the vertex-nodes incident with the exact identity $f$.
Here incidence is membership in the unordered endpoint pair, by
\noderef{Incident}.  The loopless law of \noderef{LooplessMultigraph}
makes the two entries of that pair distinct.  Thus a colored edge-node has
two distinct neighbors and cannot be a degree-one end.

Every vertex end $x$ of $R_i$ misses at least one of $a,b$.  If both were
present, \noderef{PresentColors} would provide an $a$-colored and a
$b$-colored exact edge incident with $x$.  They are distinct because
$a\ne b$, and \noderef{TwoColorIncidenceGraph} would make their two
edge-nodes distinct neighbors of $x$, contradicting degree one.  By the
complement definition \noderef{MissingColors}, $x$ therefore misses $a$
or misses $b$.

We also treat components with no edge-node.  Such a component meeting
$X$ is the singleton consisting of one vertex-node $x$ with $x\in X$:
the incidence graph is bipartite by \noderef{TwoColorIncidenceGraph}, so
there is no other kind of node to which it could be connected.  An edge
of color $a$ or $b$ incident with $x$ would give an adjacent edge-node in
the same component.  Hence \noderef{PresentColors} and
\noderef{MissingColors} show that both $a$ and $b$ are missing at $x$.
There cannot be two such singleton components, since their disjointness
would give distinct vertices of $X$ missing $a$, contrary to elementarity
in \noderef{ColoredVertexSetElementary}.

We now select two outside ends, splitting according to the status of $b$
on $T$, as required by \noderef{TreeMissingColors}.  First suppose
$b\notin\overline\varphi(T)$.  No $C_i$ is a singleton, since its tree
vertex would miss $b$.  The preceding analysis gives two disjoint
nontrivial paths and hence four distinct vertex ends.  A path end in $X$
cannot miss $b$, so it misses $a$.  Elementarity permits at most one such
tree vertex.  At least three of the four ends are therefore outside $X$;
choose any two of them.

Next suppose $b\in\overline\varphi(T)$.  If both $C_i$ are nontrivial,
each of their four distinct ends misses $a$ or $b$.  Elementarity permits
at most one end in $X$ missing $a$ and at most one end in $X$ missing
$b$.  Thus at least two ends are outside $X$.  If, after exchanging the
indices, $C_1$ is the singleton at $x\in X$, then $x$ misses both colors.
The component $C_2$ is nontrivial by the preceding singleton argument,
and neither end of its spanning path can lie in $X$: either missing color
at such an end would be repeated at the distinct vertex $x$.  Its two
ends are therefore outside $X$.  These cases produce two distinct
outside vertex ends of nontrivial component paths in every event.

For each selected outside end $a_i$, orient its spanning path toward the
component's first vertex-node belonging to $X$, and call that first vertex
$b_i$.  Such a vertex exists because the full component meets $X$.  More
explicitly, among the finitely many positions of the oriented simple path
whose node is $\operatorname{inl}(v)$ with $v\in X$, choose the least
position, and restrict the walk to the prefix ending there.  Call the
result $P_i$.  Its initial vertex is outside $X$ and misses $a$ or $b$;
its final vertex $b_i$ belongs to $X$; and no earlier vertex-node on the
prefix belongs to $X$.  The endpoints are both vertex-nodes and are
distinct, while \noderef{TwoColorIncidenceGraph} is bipartite, so the
prefix has incidence length at least two.  A prefix of a simple path is
simple.  These facts are precisely every clause of
\noderef{TreeExitPath}, so $P_i$ is a $T$-exit path.

If the selected ends came from different $C_i$, their $P_i$ are support-
disjoint because the full components are disjoint.  Their tree ends are
then distinct as well, since a common vertex-node would lie in both
supports.  It remains to check the only other possibility: the chosen
ends are the opposite ends of one spanning path.  View that simple path
as an ordered list from the first chosen end to the second.  The support
of the first-contact prefix from the left occupies the positions through
the first encounter with $X$, while the support of the reversed
first-contact prefix from the right occupies the positions from the last
encounter with $X$ to the right end.  They can overlap only if the first
and last encounters coincide, say at $b_1=b_2$.

In that exceptional situation, let $f_i$ be the exact edge-node immediately
before $\operatorname{inl}(b_i)$ on $P_i$.  Simplicity of the full path
makes $f_1\ne f_2$.  The preceding vertex-node on each $P_i$ is outside
$X$, by first contact, whereas $b_i\in X$.  The two incidences given by
\noderef{TwoColorIncidenceGraph}, interpreted by \noderef{Incident} and
the two-entry endpoint description in \noderef{LooplessMultigraph}, show
that $f_i$ has one end in $X$ and one outside.  Hence $f_i$ is an exact
member of \noderef{BoundaryEdges}.  Its incidence-graph adjacency also
says that its color is $a$ or $b$.  The hypothesis
\noderef{TreeColorSetClosed}, expanded through
\noderef{BoundaryColorEdges}, excludes color $a$ for every such exact
boundary identity.  Thus both $f_1$ and $f_2$ have color $b$.  They are
distinct exact edges incident with their common vertex, so they are
adjacent by \noderef{AdjacentEdges}; this contradicts properness in
\noderef{IsProperEdgeColoring}.  Therefore the two first-contact portions
are support-disjoint and $b_1\ne b_2$ also in this case.

The same last-edge argument applies without the exceptional equality and
shows, for each $i$, that the final exact edge-node $f_i$ of $P_i$ lies on
$P_i$, is incident with $b_i$, belongs to
$\partial_{\varphi,b}(X)$ as defined by
\noderef{BoundaryColorEdges}, and retains its individual edge identity.
Finally, the selected full-path ends $a_1,a_2$ are distinct, and each is
the outside end of its corresponding $P_i$.  All asserted witness fields
now follow.
\end{proof}
